% 解析目标 q = (100 theta_1^2 + theta_2^2)/2；轨迹按闭式递推绘制。
% 左图 eta=.015；右图为同一曲面的俯视投影，eta=.02 恰为稳定区间边界。
\begin{tikzpicture}
  \begin{axis}[name=valley,width=98mm,height=88mm,view={55}{28},
    xmin=-.4,xmax=.4,ymin=-1.2,ymax=1.2,zmin=0,zmax=9,
    xlabel={$\theta_1$},ylabel={$\theta_2$},zlabel={$q$},
    xtick={-.4,0,.4},ytick={-1,0,1},ztick={0,4,8},
    tick label style={font=\figurefont\footnotesize},
    label style={font=\figurefont\footnotesize},
    title={\figurefont （a）稳定步长：缓慢推进},
    domain=-.4:.4,y domain=-1.2:1.2,samples=17,samples y=17]
    \addplot3[mesh,draw=NNLine!65,line width=.4pt] {50*x^2+.5*y^2};
    \addplot3[NNInput,line width=1.5pt,mark=*,mark size=1.7pt,
      domain=0:8,samples=9,samples y=0,variable=\k]
      ({.35*(-.5)^\k},{.985^\k},{50*(.35*(-.5)^\k)^2+.5*(.985^\k)^2});
    \addplot3[NNInput,line width=1.3pt,domain=8:80,samples=73,samples y=0,variable=\k]
      ({.35*(-.5)^\k},{.985^\k},{50*(.35*(-.5)^\k)^2+.5*(.985^\k)^2});
    \addplot3[NNOutput,only marks,mark=*,mark size=2.3pt] coordinates {(0,0,0)};
    \node[font=\figurefont\footnotesize,text=NNInput,anchor=south west]
      at (axis cs:.35,1,6.625) {0};
    \addplot3[NNInput,only marks,mark=o,mark size=2pt]
      coordinates {(0,.2984684566,.0445417101)};
    \node[font=\figurefont\footnotesize,text=NNInput,anchor=north east]
      at (axis cs:0,.2984684566,.0445417101) {80};
  \end{axis}
  \begin{axis}[at={(valley.south east)},anchor=south west,xshift=4mm,yshift=2mm,
    width=67mm,height=88mm,axis lines=left,
    xmin=-.45,xmax=.45,ymin=-.06,ymax=1.08,
    xlabel={$\theta_1$},ylabel={$\theta_2$},xtick={-.35,0,.35},ytick={0,.5,1},
    tick label style={font=\figurefont\footnotesize},
    label style={font=\figurefont\footnotesize},
    ylabel style={at={(axis description cs:0,1)},anchor=south,rotate=0,yshift=2mm},
    title={\figurefont （b）临界步长：俯视震荡}]
    \foreach \level in {.125,.5,2,6.125}{
      \addplot[NNLine!65,line width=.5pt,domain=0:360,samples=121]
        ({sqrt(2*\level/100)*cos(x)},{sqrt(2*\level)*sin(x)});
    }
    \addplot[NNLine,densely dashed,line width=.5pt]
      coordinates {(0,0) (0,1.04)};
    \addplot[NNGradient,line width=1pt,domain=0:20,samples=21,variable=\k]
      ({.35*(-1)^\k},{.98^\k});
    \addplot[NNGradient,only marks,mark=*,mark size=1.5pt,
      domain=0:6,samples=7,variable=\k]
      ({.35*(-1)^\k},{.98^\k});
    \addplot[NNGradient,only marks,mark=o,mark size=2pt]
      coordinates {(.35,.6676079718)};
    \node[font=\figurefont\footnotesize,text=NNGradient,anchor=south east]
      at (axis cs:.35,1) {0};
    \node[font=\figurefont\footnotesize,text=NNGradient,anchor=north east]
      at (axis cs:.35,.6676079718) {20};
    \draw[NNGradient,line width=1pt,-{Stealth[length=1.8mm]}]
      (axis cs:-.12,.935007024) -- (axis cs:.12,.928553136);
    \draw[NNGradient,line width=1pt,-{Stealth[length=1.8mm]}]
      (axis cs:.12,.7486894373) -- (axis cs:-.12,.7435216068);
    \addplot[NNGradient,only marks,mark=square*,mark size=2pt]
      coordinates {(-.35,0) (.35,0)};
    \addplot[NNOutput,only marks,mark=*,mark size=2.3pt] coordinates {(0,0)};
  \end{axis}
\end{tikzpicture}
